\documentclass[11pt,a4paper]{article}

% --- Encoding ---
\usepackage[utf8]{inputenc}
\usepackage[T1]{fontenc}
\usepackage[spanish,es-noshorthands]{babel}
\usepackage{lmodern}

% --- Layout ---
\usepackage[margin=2.5cm]{geometry}
\usepackage{parskip}
\setlength{\parindent}{0pt}
\setlength{\parskip}{0.5em plus 0.2em minus 0.1em}

% --- Tables and graphics ---
\usepackage{booktabs}
\usepackage{array}
\usepackage{enumitem}
\usepackage{xcolor}
\usepackage{graphicx}
\usepackage{tikz}
\usetikzlibrary{positioning,shapes.geometric,arrows.meta,fit,backgrounds,calc,decorations.pathreplacing}
\usepackage{caption}
\usepackage{amsmath,amssymb}

% --- Code listings ---
\usepackage{listings}
\definecolor{codegray}{rgb}{0.96,0.96,0.96}
\definecolor{codeborder}{rgb}{0.85,0.85,0.85}
\definecolor{codecomment}{rgb}{0.30,0.55,0.30}
\definecolor{codekeyword}{rgb}{0.10,0.30,0.70}
\definecolor{codestring}{rgb}{0.65,0.15,0.10}

\lstdefinestyle{cstyle}{
    language=C,
    backgroundcolor=\color{codegray},
    basicstyle=\ttfamily\footnotesize,
    breaklines=true,
    frame=single,
    rulecolor=\color{codeborder},
    showstringspaces=false,
    keywordstyle=\color{codekeyword}\bfseries,
    commentstyle=\color{codecomment}\itshape,
    stringstyle=\color{codestring},
    aboveskip=1em,
    belowskip=1em,
}

\lstdefinestyle{textstyle}{
    backgroundcolor=\color{codegray},
    basicstyle=\ttfamily\footnotesize,
    breaklines=true,
    frame=single,
    rulecolor=\color{codeborder},
    showstringspaces=false,
    aboveskip=1em,
    belowskip=1em,
    columns=fullflexible,
}

\lstset{style=textstyle}

% --- Definition / Idea / Important boxes ---
\usepackage[most]{tcolorbox}
\newtcolorbox{definicion}[1][]{
    colback=blue!5!white,
    colframe=blue!50!black,
    fonttitle=\bfseries,
    title=Definición,
    sharp corners=southwest,
    rounded corners=northwest,
    arc=2pt,
    boxrule=0.6pt,
    left=8pt,right=8pt,top=4pt,bottom=4pt,
    #1
}
\newtcolorbox{idea}[1][]{
    colback=orange!5!white,
    colframe=orange!60!black,
    fonttitle=\bfseries,
    title=Idea intuitiva,
    sharp corners=southwest,
    rounded corners=northwest,
    arc=2pt,
    boxrule=0.6pt,
    left=8pt,right=8pt,top=4pt,bottom=4pt,
    #1
}
\newtcolorbox{importante}[1][]{
    colback=red!5!white,
    colframe=red!50!black,
    fonttitle=\bfseries,
    title=Importante,
    sharp corners=southwest,
    rounded corners=northwest,
    arc=2pt,
    boxrule=0.6pt,
    left=8pt,right=8pt,top=4pt,bottom=4pt,
    #1
}

% --- Hyperlinks ---
\usepackage[colorlinks=true,
            linkcolor=blue!60!black,
            urlcolor=blue!50!black,
            citecolor=blue!60!black]{hyperref}

% --- Memory diagram color palette ---
\definecolor{mem1}{HTML}{4C72B0}
\definecolor{mem2}{HTML}{DD8452}
\definecolor{mem3}{HTML}{55A868}
\definecolor{mem4}{HTML}{8172B2}
\definecolor{memused}{HTML}{C44E52}
\definecolor{memfree}{HTML}{8C8C8C}

% --- Title ---
\title{
    \vspace{-2.5em}
    \textbf{Resumen teórico --- Sistemas Operativos} \\[0.3em]
    \large Notas del curso: \emph{Gestión de la memoria}
}
\author{
    Tomás Rodeghiero \\
    \small Sistemas Operativos --- Universidad Nacional de Río Cuarto
}
\date{2026}

\begin{document}

\maketitle

\begin{abstract}
\noindent
Este documento es un resumen teórico del capítulo \emph{Gestión de la
memoria} de las \emph{Notas del curso} de Sistemas Operativos (Marcelo
Arroyo, UNRC). Se presentan los conceptos básicos sobre la memoria
principal y la jerarquía de memorias, las estrategias de asignación de
memoria desde el heap (bloques de tamaño variable, bloques de tamaño
fijo y \emph{buddy system}), los mecanismos de protección por registros
base/límite y segmentación, el código independiente de la posición, el
esquema de paginado con sus tablas de páginas multinivel y, finalmente,
los \emph{Translation Lookaside Buffers} (TLB) como caché de
traducciones. El objetivo es servir de material de estudio para preparar
los prácticos asociados a la gestión de memoria y el segundo parcial.
\end{abstract}

\tableofcontents
\bigskip

% =====================================================================
\section{Introducción}
% =====================================================================

La \textbf{memoria principal} o RAM (\emph{Random Access Memory}) es,
desde el punto de vista del sistema, un arreglo de bytes en el que cada
byte se identifica por su \emph{dirección} o posición. Es una memoria
\textbf{volátil}: el contenido se pierde cuando se corta la energía.
La CPU interactúa directamente con la RAM en cada ciclo de ejecución
de cada instrucción.

\begin{idea}
En las arquitecturas \emph{Von Neumann} las instrucciones del programa
y los datos sobre los que opera \emph{comparten} la misma memoria, a
diferencia de la arquitectura Harvard que separa físicamente ambos
espacios.
\end{idea}

En un sistema multitarea la memoria física debe ser
\textbf{particionada lógicamente} para separar las instrucciones y los
datos de cada programa en ejecución. Esto le impone al SO dos
responsabilidades inseparables:

\begin{itemize}[leftmargin=*]
    \item \textbf{Asignar} memoria al kernel y a cada uno de los
        procesos: decidir qué regiones de la RAM le tocan a quién, y
        cómo se liberan cuando ya no se usan.
    \item \textbf{Proteger} los espacios de memoria: garantizar que un
        proceso se ejecute \emph{confinado}, es decir, que sólo pueda
        acceder a sus propias regiones y no pueda invadir las de otros
        procesos ni las del kernel.
\end{itemize}

\subsection{El ciclo de instrucción}

La CPU ejecuta repetidamente un ciclo con tres pasos:

\begin{enumerate}[leftmargin=*]
    \item \textbf{Fetch (lectura) de la instrucción.} Se carga en un
        registro interno la instrucción cuya dirección está en el
        \emph{program counter} (\texttt{pc}).
    \item \textbf{Decodificación.} Se identifica el tipo de instrucción
        y sus operandos.
    \item \textbf{Ejecución} de la instrucción, que en general puede
        requerir \emph{leer} o \emph{escribir} en memoria.
\end{enumerate}

Como se ve, todo paso del ciclo toca memoria; por eso la velocidad y la
correcta gestión de la RAM determinan en gran medida el rendimiento del
sistema.

% =====================================================================
\section{Jerarquía de memorias}
% =====================================================================

Las distintas memorias del sistema no son equivalentes: hay un
\emph{trade-off} clásico entre \textbf{velocidad} y \textbf{capacidad}
(y costo). La CPU emplea \emph{registros} internos para ejecutar las
instrucciones; son las memorias más rápidas, pero hay muy pocos. Cada
operación de lectura o escritura hacia/desde la RAM puede tomar varios
ciclos de reloj, lo que obliga a la CPU a \emph{frenarse} (\emph{stall})
hasta que el dato esté disponible.

Para atenuar esta diferencia de velocidades se introducen las memorias
\textbf{caché}: memorias de menor capacidad pero de mucha mayor velocidad
de acceso que la RAM convencional, ubicadas entre los registros y la
RAM.

\begin{definicion}
Una memoria \textbf{caché} está diseñada para el acceso a los datos
\emph{por contenido} en lugar de por dirección. Cuando la CPU referencia
una dirección de memoria:
\begin{itemize}[leftmargin=*,nosep]
    \item si el contenido de esa dirección ya está en la caché
        (\emph{hit}), se recupera desde allí;
    \item si no se encuentra (\emph{miss}), se accede a la RAM y se
        almacena en la caché para los accesos futuros.
\end{itemize}
\end{definicion}

Finalmente, para almacenar programas y datos en forma \textbf{permanente}
se usan dispositivos secundarios como discos magnéticos, unidades de
estado sólido u otros, que son los más lentos pero también los de mayor
capacidad. La arquitectura define entonces una \emph{jerarquía} de
memorias ordenadas por velocidad de acceso (de mayor a menor):

\begin{center}
\begin{tikzpicture}[
    every node/.style={font=\ttfamily\small},
    box/.style={draw,rounded corners=2pt,minimum width=2.1cm,minimum height=1cm,align=center},
    arrow/.style={-Latex,thick}
]
\node[box,fill=mem1!20] (reg) {cpu\\registers};
\node[box,fill=mem3!20,right=0.6cm of reg] (cache) {cache};
\node[box,fill=mem2!20,right=0.6cm of cache] (ram) {main\\memory};
\node[box,fill=memfree!20,right=0.6cm of ram] (disc) {disc};

\draw[arrow] (cache) -- (reg);
\draw[arrow] (ram) -- (cache);
\draw[arrow] (disc) -- (ram);
\end{tikzpicture}
\end{center}

\begin{idea}
A medida que se asciende en la jerarquía (a la izquierda) la memoria es
más rápida, más cara y de menor capacidad. La idea detrás de la caché es
explotar la \emph{localidad de referencia}: los programas tienden a
acceder repetidamente a un conjunto pequeño de direcciones cercanas en
el tiempo.
\end{idea}

% =====================================================================
\section{Asignación (\emph{allocation}) de memoria}
% =====================================================================

El kernel debe \emph{asignar} y \emph{liberar} bloques de memoria, tanto
para los procesos de usuario como para sus propias estructuras de datos
internas.

Una vez que el kernel se cargó en memoria durante el proceso de arranque
(\emph{boot}), reconoce el área de memoria libre disponible para
asignar. Esta área se conoce como \textbf{heap}.

Sobre el heap se aplican dos grandes estrategias clásicas:

\begin{enumerate}[leftmargin=*]
    \item \textbf{Asignación de bloques de tamaño variable:} el heap se
        particiona lógicamente y se asignan bloques cuyo tamaño se
        ajusta al pedido.
    \item \textbf{Asignación de bloques de tamaño fijo:} el heap se
        particiona en bloques todos del mismo tamaño y se asignan
        bloques completos (o varios contiguos) por cada pedido.
\end{enumerate}

\begin{center}
\begin{tikzpicture}[
    every node/.style={font=\ttfamily\small},
    block/.style={draw,minimum width=3.2cm,minimum height=0.8cm,align=center},
]
% Variable size blocks
\node[font=\bfseries\small] at (0,4.6) {Variable size blocks};
\node[block,fill=memused!30] (v1) at (0,3.7) {used size=n1};
\node[block,fill=memfree!20,below=0pt of v1] (v2) {free size=n2};
\node[block,fill=memused!30,below=0pt of v2,minimum height=1.0cm] (v3) {used size=n3};
\node[block,fill=memfree!20,below=0pt of v3] (v4) {free size=n4};
\node[block,below=0pt of v4,minimum height=0.5cm] (vdots) {\ldots};

% Fixed size blocks
\node[font=\bfseries\small] at (5.5,4.6) {Fixed size blocks};
\node[block,fill=memused!30] (f1) at (5.5,3.95) {size=n};
\node[block,below=0pt of f1,fill=memfree!10] (f2) {};
\node[block,below=0pt of f2,fill=memused!30] (f3) {};
\node[block,below=0pt of f3,fill=memfree!10] (f4) {};
\node[block,below=0pt of f4,minimum height=0.5cm] (fdots) {\ldots};
\end{tikzpicture}
\end{center}

% ---------------------------------------------------------------------
\subsection{Heap con bloques de tamaño variable}
% ---------------------------------------------------------------------

La primera estrategia se implementa habitualmente mediante una
\textbf{lista de bloques libres} (\emph{freelist}):

\begin{itemize}[leftmargin=*]
    \item Inicialmente existe un único bloque libre del tamaño total
        del heap.
    \item Ante una solicitud de un bloque de tamaño $n$, se busca en la
        lista un bloque libre $b$ con $b.\textit{size} \geq n$. Si
        $b.\textit{size} > n$, se \emph{particiona}: se asigna un bloque
        de tamaño $n$ y queda un nuevo bloque libre $b'$ con
        $b'.\textit{size} = b.\textit{size} - n$. Si no existe ningún
        bloque que satisfaga el pedido, se retorna error.
    \item Ante la liberación de un bloque, se inserta nuevamente en la
        lista de libres.
\end{itemize}

El algoritmo concreto que se use para buscar el bloque adecuado
afecta directamente al rendimiento del sistema. Las tres estrategias
clásicas son:

\begin{enumerate}[leftmargin=*]
    \item \textbf{First fit.} Se asigna el primer bloque encontrado con
        tamaño $\geq n$. Es la más rápida porque corta la búsqueda en
        cuanto encuentra un candidato.
    \item \textbf{Best fit.} Se asigna el bloque que \emph{mejor ajusta},
        es decir, un bloque $b$ con
        $b.\textit{size} = m$, $m \geq n$ y para todo $b' \in
        \textit{freelist}$ se cumple
        $b'.\textit{size} \geq b.\textit{size}$. Minimiza el desperdicio
        por pedido, pero suele dejar fragmentos pequeños y exige
        recorrer toda la lista.
    \item \textbf{Worst fit.} Se asigna el bloque libre más grande. La
        idea es que el remanente sea lo suficientemente grande como
        para poder ser reutilizado.
\end{enumerate}

\paragraph{Fragmentación externa.}
A medida que se hacen sucesivos pedidos y liberaciones, el heap puede
ir \textbf{fragmentándose} en un gran número de bloques libres pequeños.
Aunque la suma de esos bloques pueda ser grande, ningún bloque
\emph{individual} alcanza para satisfacer un pedido grande. Este
problema se conoce como \textbf{fragmentación externa}.

\begin{importante}
La fragmentación externa es el costo principal del esquema de tamaño
variable: el heap puede tener mucha memoria libre y aun así rechazar un
pedido por no poder ofrecer un bloque contiguo suficientemente grande.
\end{importante}

\paragraph{Semi-compactación.}
Una estrategia que alivia la fragmentación es: ante la liberación de un
bloque, si quedan bloques libres \emph{adyacentes}, fusionarlos o
\emph{compactarlos} en un único bloque libre de mayor tamaño. Esta
técnica se conoce como \textbf{semi-compactación}.

% ---------------------------------------------------------------------
\subsection{Heap con bloques de tamaño fijo}
% ---------------------------------------------------------------------

En esta variante el heap se particiona en $m$ bloques de un tamaño fijo
$s$.

Ante un requerimiento de un bloque de tamaño $n$:

\begin{itemize}[leftmargin=*,nosep]
    \item Si $n \leq s$, se asigna cualquier bloque libre.
    \item Si $n > s$, se debe asignar el mínimo número $k$ de bloques
        \emph{contiguos} tales que $k \times s \geq n$.
\end{itemize}

Una ventaja muy importante en términos de uso de memoria es que las
listas (o conjuntos) de bloques libres y usados pueden representarse
como un \textbf{bit vector}: un bit por bloque, encendido si está
asignado y apagado si está libre.

\paragraph{La elección del tamaño \texorpdfstring{$s$}{s}.}
El tamaño $s$ es \emph{crítico}:

\begin{itemize}[leftmargin=*]
    \item Si $s$ es grande, se agrava la \textbf{fragmentación interna}:
        cada pedido recibe un bloque completo, pero el remanente
        $s - n$ no se aprovecha.
    \item Si $s$ es pequeño, una asignación grande requiere encontrar
        varios bloques libres \emph{contiguos}; si esa contigüidad no
        existe, el pedido no se puede satisfacer aunque haya muchos
        bloques libres dispersos.
\end{itemize}

\begin{idea}
La fragmentación interna es ``desperdicio dentro del bloque asignado''
(memoria que el proceso recibió pero no usa). La fragmentación externa
es ``desperdicio entre bloques'' (memoria libre que no se puede
asignar porque no es contigua). Ambos esquemas tienen sus
contraindicaciones.
\end{idea}

% ---------------------------------------------------------------------
\subsection{El sistema compañero (\emph{buddy system})}
% ---------------------------------------------------------------------

Para combinar las ventajas de ambas estrategias se diseñan soluciones
\emph{híbridas}. El \textbf{buddy system} es el ejemplo clásico.

La idea es particionar el heap en bloques cuyos tamaños sean
\emph{potencias de 2}, desde un mínimo $2^l$ hasta un máximo $2^h$. Por
ejemplo, con $l = 16$ y $h = 20$ los bloques posibles son de 64\,KB,
128\,KB, 256\,KB, 512\,KB y 1024\,KB.

\paragraph{Algoritmo de asignación.}
\begin{enumerate}[leftmargin=*]
    \item Inicialmente, el heap se particiona en bloques del tamaño
        máximo $2^h$ (1024\,KB en el ejemplo).
    \item Ante un requerimiento de tamaño $n$:
        \begin{enumerate}[leftmargin=*,label=\arabic*.]
            \item Seleccionar un bloque libre $b$ con mejor ajuste.
            \item Si $b.\textit{size} \geq 2n$, dividir $b$ en dos
                bloques $b_0$ y $b_1$ de tamaño $b.\textit{size}/2$.
                Tomar $b = b_0$ y repetir este paso mientras
                $b.\textit{size} \geq 2n$ (y siempre que
                $b.\textit{size} > 2^l$).
            \item Asignar el bloque $b$.
        \end{enumerate}
\end{enumerate}

\paragraph{Liberación.}
Cuando se libera un bloque se aplica el algoritmo a la inversa: si el
bloque \emph{compañero} (\emph{buddy}) contiguo es del mismo tamaño y
también está libre, ambos se unen en un único bloque del doble de
tamaño, mientras esto se pueda repetir y sin pasar de $2^h$.

\begin{importante}
El buddy system es un compromiso entre fragmentación interna y externa:
mantiene bloques de tamaño potencia de 2 (lo que simplifica la
contabilidad y la fusión), pero al mismo tiempo permite que el heap se
particione dinámicamente según los pedidos. Es la base de muchos
\emph{allocators} de kernel reales (por ejemplo el de Linux).
\end{importante}

% =====================================================================
\section{Protección}
% =====================================================================

Asignar memoria es sólo la mitad del problema: el SO también debe
garantizar que cada proceso \textbf{sólo pueda acceder a su área de
memoria}. Supongamos que un proceso recibe su memoria de forma
\emph{contigua}; entonces basta con que el hardware pueda
\emph{detectar} todo acceso fuera de rango.

\subsection{Registros base y límite}

El mecanismo más simple consiste en un registro de control que contiene
la \textbf{dirección base} y la \textbf{dirección límite} del área de
memoria asignada al proceso. Además puede incluir \textbf{flags} de
\emph{permisos de acceso}: lectura, escritura, ejecución, etc.

\begin{center}
\begin{tikzpicture}[
    every node/.style={font=\ttfamily\small},
    reg/.style={draw,minimum height=0.9cm,minimum width=1.4cm,align=center},
    mem/.style={draw,minimum width=2.6cm,align=center}
]
\node[reg] (flags) {flags};
\node[reg,right=0pt of flags] (base)  {base};
\node[reg,right=0pt of base]  (limit) {limit};

\node[mem,right=2.0cm of limit,fill=mem3!15,minimum height=2.6cm] (mem)
    {\\\\process\\memory\\\\};
\node[anchor=south] at (mem.north) {\ttfamily ...};
\node[anchor=north] at (mem.south) {\ttfamily ...};

\draw[-Latex,thick] (base.east) .. controls +(1.0,0.7) and +(-1.0,0.5) .. ($(mem.north west)+(0,-0.3)$);
\draw[-Latex,thick] (limit.east) .. controls +(1.0,-0.8) and +(-1.0,-0.5) .. ($(mem.south west)+(0,0.3)$);
\end{tikzpicture}
\end{center}

Cuando el SO realiza un \emph{context switch} a un proceso, carga las
direcciones correspondientes en estos registros. La CPU, en cada acceso
a memoria, verifica que la dirección referenciada caiga en el rango
permitido:
\[
\textit{base} \leq \textit{address} \leq \textit{limit}.
\]

\paragraph{Memory fault exception.}
Si una instrucción referencia una dirección fuera de los límites, o si
la operación de lectura/escritura no respeta los permisos definidos en
los flags, la CPU genera una \textbf{memory fault exception}. El SO
atrapará la excepción con algún \emph{exception handler} y, si no es
recuperable, eliminará al proceso. Este error es exactamente el
conocido \textbf{segmentation fault} de los sistemas tipo UNIX.

\subsection{Segmentación}

El esquema de un único par base/límite es muy rígido: obliga a que toda
la memoria del proceso esté en un bloque contiguo, lo que no siempre se
puede lograr en un sistema fragmentado.

La \textbf{segmentación} extiende la idea: se utilizan \emph{varios}
pares base/límite, uno por cada \emph{segmento} lógico del proceso
(\emph{texto}, \emph{datos}, \emph{stack}, \emph{heap}, etc.). Así cada
segmento puede vivir en una región de memoria distinta y los segmentos
no necesitan ser físicamente contiguos entre sí.

\begin{idea}
A modo de ejemplo, la arquitectura Intel IA-32 (x86) soporta varios
registros de segmento (base + límite):
\textbf{CS} (\emph{code segment}), \textbf{DS} (\emph{data segment}),
\textbf{SS} (\emph{stack segment}) y otros como \textbf{ES}, \textbf{GS},
\textbf{FS}.
\end{idea}

% =====================================================================
\section{Código independiente de la posición (PIC)}
% =====================================================================

Con la maquinaria de protección por segmentos aparece otro beneficio:
el código de un programa puede ejecutarse \emph{independientemente} de
la dirección base donde se cargó. Esto se logra con
\textbf{código independiente de la posición} (\emph{position-independent
code}, \textbf{PIC}).

\begin{definicion}
En el código PIC, las direcciones de memoria que aparecen como
operandos de las instrucciones no son \emph{absolutas} sino
\textbf{relativas} a una dirección base contenida en un registro de
segmento.
\end{definicion}

Los registros de segmento contienen las direcciones base; la CPU
\emph{transforma} cada dirección relativa en una \textbf{dirección
física} sumándole la dirección base correspondiente al segmento.

\paragraph{Notación en x86.}
En x86, una dirección de memoria en una instrucción se denota como
\texttt{segment:offset}. Algunos casos:

\begin{itemize}[leftmargin=*,nosep]
    \item Un salto \verb|jmp address| es equivalente a
        \verb|jmp cs:address|.
    \item Una instrucción de lectura/escritura desde/hacia memoria
        como \verb|mov %eax, address| es equivalente a
        \verb|mov %eax, ds:address|.
    \item Las instrucciones sobre la pila (\verb|push| / \verb|pop|)
        refieren por omisión al segmento \verb|ss|.
\end{itemize}

\begin{idea}
PIC es lo que permite, por ejemplo, que el mismo binario de una
biblioteca compartida (\emph{shared library}) se pueda cargar en una
dirección distinta en cada proceso sin tener que reescribir su código.
También es la base del esquema ASLR moderno (\emph{Address Space Layout
Randomization}).
\end{idea}

% =====================================================================
\section{Paginado}
% =====================================================================

La segmentación resuelve el problema de no necesitar bloques contiguos
grandes para todo el proceso, pero todavía sufre de fragmentación
externa al manejar segmentos enteros. La solución estándar de los SO
modernos es el \textbf{paginado}.

\begin{definicion}
El \textbf{paginado} consiste en particionar la memoria física en
bloques pequeños del \emph{mismo tamaño}, llamados \textbf{páginas} o
\emph{frames}, típicamente de 4\,KB. Cada espacio de direcciones lógicas
de un proceso (y del kernel) se representa como un conjunto de páginas
que pueden vivir en cualquier ubicación física.
\end{definicion}

\begin{center}
\begin{tikzpicture}[
    every node/.style={font=\ttfamily\small},
    page/.style={draw,minimum width=2.6cm,minimum height=0.7cm,align=center}
]
\node[page,fill=mem1!15] (p0) {page 0};
\node[page,fill=mem2!15,below=0pt of p0] (p1) {page 1};
\node[page,below=0pt of p1] (pd) {\ldots};
\node[page,fill=mem3!15,below=0pt of pd] (pn) {page n-1};

\node[anchor=west,xshift=0.2cm] at (p0.east) {0};
\node[anchor=west,xshift=0.2cm] at ($(p0.south east)!0.5!(p1.north east)$) {4096};
\node[anchor=west,xshift=0.2cm] at ($(p1.south east)!0.5!(pd.north east)$) {8192};
\node[anchor=west,xshift=0.2cm] at (pn.east) {max address};
\end{tikzpicture}
\end{center}

Con paginado es posible definir un \textbf{espacio de memoria} para cada
proceso (y el kernel) como el conjunto de páginas a las que puede
acceder. Estos conjuntos se representan en memoria como arreglos y se
conocen como \textbf{tablas de páginas} (\emph{page tables}).

\paragraph{Direcciones lógicas vs físicas.}
Las páginas asignadas a un proceso no necesitan estar
\emph{físicamente contiguas}, aunque para el proceso aparecen como si lo
estuvieran. Esto requiere de un mecanismo de \emph{traducción} de
direcciones \textbf{lógicas o virtuales} a direcciones
\textbf{físicas o reales}. Formalmente, se necesita una función
\[
\textit{map} : \textit{va} \to \textit{pa},
\]
donde $\textit{va}$ es una dirección lógica (\emph{virtual address}) y
$\textit{pa}$ es la dirección física correspondiente.

El hardware con soporte de paginado implementa esta función mediante una
\emph{tabla de páginas} que representa la traducción
$\textit{paging}: \textit{va} \to \textit{pa}$. Cada entrada de la
tabla se llama \textbf{page table entry} (\textbf{pte}) y básicamente
es un par
$(\textit{ppn},\;\textit{flags})$,
donde:

\begin{itemize}[leftmargin=*,nosep]
    \item \texttt{ppn} (\emph{physical page number}) es el número de
        página física correspondiente (las páginas se numeran desde 0);
    \item \texttt{flags} son permisos de acceso y bits de control.
\end{itemize}

Por cada proceso y para el kernel, el SO construye su propio espacio de
memoria asociando una tabla de páginas distinta.

\subsection{Traducción de una dirección lógica}

\begin{importante}
En paginado, una dirección que aparece en una instrucción del programa
es, en realidad, una \emph{dirección lógica} y generalmente se
interpreta como un par
$(\textit{index},\;\textit{offset})$.
\end{importante}

\begin{itemize}[leftmargin=*,nosep]
    \item El \textbf{index} (\emph{virtual page number}, \texttt{vpn})
        identifica la entrada de la tabla de páginas asociada al
        proceso.
    \item El \textbf{offset} es el desplazamiento \emph{dentro} de la
        página.
\end{itemize}

La dirección física se obtiene como
\[
\textit{pa} = \textit{pagetable}[\textit{index}].\textit{ppn} \times PS
              + \textit{offset},
\]
donde $PS$ es el tamaño de página (\emph{page size}).

\begin{center}
\begin{tikzpicture}[
    every node/.style={font=\ttfamily\small},
    cell/.style={draw,minimum height=0.7cm,align=center},
]
% Logical address
\node[font=\bfseries\small] at (0,2.3) {logical address};
\node[cell,fill=mem1!15,minimum width=1.4cm] (idx) at (-0.8,1.6) {index};
\node[cell,fill=mem2!15,minimum width=1.4cm] (off) at (0.6,1.6) {offset};

% Pagetable
\node[font=\bfseries\small] at (3.6,2.3) {pagetable};
\node[cell,minimum width=1.0cm] (p0a) at (3.0,1.6) {ppn};
\node[cell,minimum width=1.0cm] (p0b) at (4.1,1.6) {flags};
\node[cell,minimum width=1.0cm,fill=mem3!20] (p1a) at (3.0,0.9) {ppn};
\node[cell,minimum width=1.0cm,fill=mem3!20] (p1b) at (4.1,0.9) {flags};
\node[anchor=east] at (p0a.west) {0};
\node[anchor=east] at (p1a.west) {1};
\node at (3.6,0.15) {\ldots};

% Frame
\node[cell,minimum width=2.0cm,minimum height=1.3cm,fill=mem4!20] (frame) at (7.5,1.0) {page\\(frame)};

% Phys addr
\node[font=\bfseries\small] at (10.7,1.6) {physical};
\node[font=\bfseries\small] at (10.7,1.1) {address};

\draw[-Latex,thick] (idx.south) -- ++(0,-0.4) -| (p1a.west);
\draw[-Latex,thick] (p1b.east) -- (frame.west);
\draw[-Latex,thick] (off.east) -- ++(0.4,0) -- ++(0,-0.6) -| ($(frame.south)+(0.3,0)$);
\draw[-Latex,thick] (frame.east) -- ++(0.7,0);
\end{tikzpicture}
\end{center}

\paragraph{Tamaño del offset.}
El offset contiene tantos bits como sea necesario según el tamaño de
página. Por ejemplo, si el tamaño de página es 4096\,bytes
($2^{12}$), el offset necesita 12\,bits.

\paragraph{Page fault.}
La CPU genera una excepción de tipo \textbf{page fault} si una dirección
lógica indexa en una \emph{pte} sin permisos de acceso o inválida (por
ejemplo, una entrada cuyo bit \texttt{V} está apagado).

\subsection{Espacios de direcciones por defecto}

Con paginado es posible que cada proceso tenga, por ejemplo, un espacio
de direcciones de la forma $[0, \textit{size}]$ donde \textit{size} es
la memoria requerida por el proceso. Así todos los programas se pueden
compilar con opciones de memoria \emph{por default} y pueden coexistir
en memoria \emph{referenciando las mismas direcciones lógicas}: en
tiempo de ejecución cada una se traducirá a direcciones físicas
distintas.

\begin{idea}
Dos procesos pueden leer la dirección lógica \texttt{0x00001000} sin
pisarse: la tabla de páginas de cada uno mapea ese índice a frames
físicos distintos. El hardware se encarga, transparentemente, de que
cada proceso ``vea'' su propio espacio de memoria.
\end{idea}

\subsection{Páginas dispersas y heap allocator}

Cuando un proceso usa paginado, un bloque de memoria con páginas
\emph{lógicamente contiguas} no necesariamente se corresponde con
páginas \emph{físicamente contiguas}. Esto permite que un
\textbf{heap allocator} le asigne páginas al proceso con cualquier
dirección física, mientras quede mapeada a una dirección lógica del
proceso.

A modo de ejemplo, un proceso podría tener:

\begin{itemize}[leftmargin=*,nosep]
    \item 2 páginas (8\,KB) para el área de \emph{código} (texto);
    \item 1 página para \emph{datos};
    \item 1 página para el \emph{stack};
    \item un \emph{hueco} de páginas inválidas en el medio para detectar
        \textbf{stack overflow}.
\end{itemize}

Las páginas de código y datos resultan lógicamente contiguas, pero
físicamente pueden estar dispersas por toda la memoria.

\paragraph{Ejemplo numérico.}
Asumiendo direcciones de 32\,bits con los 20\,bits más significativos
como \emph{index} y los 12 menos significativos como \emph{offset}, en
páginas de 4\,KB:

Si el espacio de direcciones lógicas es $[0, sz = \texttt{0x5000}]$
(20\,KB), una instrucción que referencie la dirección lógica
\texttt{0x00002004} (área de datos) se interpreta con
\texttt{index = 2} y \texttt{offset = 0x004}. Si la \emph{pte} de la
entrada 2 contiene \texttt{ppn = 11}, la dirección física final es:
\[
11 \times 4096 + 4 = \texttt{0xb004}.
\]

\subsection{Hardware vs.\ memoria direccionable}

Las tablas de páginas deberían tener suficientes entradas como para
\emph{barrer} toda la memoria direccionable por el sistema (esto es
crítico en el kernel, que debe poder acceder a toda la memoria).

\begin{importante}
Es común que el hardware tenga \emph{menos RAM física} que el espacio de
direcciones lógicas posibles. Por ejemplo, en el modelo
\textbf{RISC-V RV39} la CPU es de 64\,bits (espacio de direcciones
$[0, 2^{64}]$), pero el bus físico de direcciones es de sólo 39\,bits
(500\,GB).
\end{importante}

\subsection{Páginas grandes y tablas multinivel}

Con espacios de direcciones grandes (por ejemplo 4\,GB en una
arquitectura de 32\,bits) el número de entradas en una tabla de páginas
puede ser enorme. Representar una tabla plana con tantas entradas
requeriría asignar un bloque de memoria física contigua de gran tamaño,
algo costoso o directamente inviable.

Para mitigarlo existen dos herramientas combinables:

\begin{enumerate}[leftmargin=*]
    \item \textbf{Páginas más grandes.} El hardware moderno permite
        definir tablas que refieran a páginas de mayor tamaño (4\,MB o
        más). Esto reduce la cantidad total de entradas pero a costa
        de aumentar la \emph{fragmentación interna}: cada página
        desperdicia más memoria si no se llena por completo.
    \item \textbf{Tablas multinivel.} Aún con páginas pequeñas, los
        procesos suelen dejar muchos \emph{huecos} en su espacio de
        direcciones lógicas, lo que en una tabla plana sería
        desperdicio. Por eso las tablas de páginas se representan como
        \textbf{árboles de $n$ niveles}.
\end{enumerate}

Cada nodo del árbol tiene el tamaño de una página. Sólo se reservan
los nodos correspondientes a regiones efectivamente usadas del espacio
de direcciones; las regiones vacías no consumen memoria.

\subsection{Caso de estudio: paginado en RISC-V RV-32}

En la arquitectura \textbf{RISC-V RV-32} el espacio de direcciones es
de 32\,bits (4\,GB) y la tabla de páginas tiene la forma de un
\textbf{árbol de dos niveles}. El registro \texttt{satp}
(\emph{Supervisor Address Translation and Protection}), de acceso
privilegiado, contiene el número de página física de la \emph{raíz} del
árbol.

\begin{center}
\begin{tikzpicture}[
    every node/.style={font=\ttfamily\small},
    cell/.style={draw,minimum height=0.7cm,align=center},
]
\node[font=\bfseries\small] at (3.2,2.3) {virtual address};
\node[cell,minimum width=1.4cm,fill=mem1!20] (vpn1) at (1.2,1.6) {vpn1};
\node[cell,minimum width=1.4cm,fill=mem2!20] (vpn0) at (2.7,1.6) {vpn0};
\node[cell,minimum width=2.0cm,fill=mem3!20] (off)  at (4.6,1.6) {offset};
\node[font=\scriptsize] at (1.2,2.1) {10 bits};
\node[font=\scriptsize] at (2.7,2.1) {10 bits};
\node[font=\scriptsize] at (4.6,2.1) {12 bits};
\end{tikzpicture}
\end{center}

Una dirección lógica se divide en tres partes:

\begin{itemize}[leftmargin=*,nosep]
    \item \texttt{vpn1} (10\,bits más significativos): índice o
        \emph{virtual page number} al \emph{nodo raíz} del árbol. Su
        \texttt{pte} contiene el \texttt{ppn} del nodo hijo.
    \item \texttt{vpn0} (10\,bits siguientes): índice en el nodo hoja.
        Su \texttt{pte} contiene el \texttt{ppn} de la página o
        \emph{frame} destino.
    \item \texttt{offset} (12\,bits menos significativos): desplazamiento
        dentro de la página destino.
\end{itemize}

\paragraph{Ejemplo.}
La dirección lógica \texttt{0x0000100f} se interpreta como:

\begin{center}
\begin{tabular}{@{}lll@{}}
\textbf{vpn1} & \textbf{vpn0} & \textbf{offset} \\
\midrule
\texttt{0000000011} & \texttt{0000000001} & \texttt{000000001111} \\
\end{tabular}
\end{center}

La traducción procede así:

\begin{enumerate}[leftmargin=*]
    \item Comienza en el nodo raíz en la entrada $\texttt{vpn1}=3$
        (tercera entrada).
    \item Suponiendo que esa \texttt{pte} contiene un \texttt{ppn} que
        lleva al nodo hijo (digamos $y$), se indexa en el nodo hoja
        $y$ en la entrada $\texttt{vpn0}=1$.
    \item En esa entrada el \texttt{ppn} = $k$ lleva a la dirección
        base física del frame destino.
    \item La dirección física final es la dirección base del frame
        sumada al offset (\texttt{0x00f}).
\end{enumerate}

Formalmente, la función de traducción es
\[
\textit{pa} = \bigl(\textit{satp}.\textit{ppn}[\textit{va}.\textit{vpn1}].\textit{ppn} \times PS\bigr)[\textit{va}.\textit{vpn0}].\textit{ppn} \times PS + \textit{va}.\textit{offset}.
\]

\subsection{Formato de una \emph{page table entry} en RV-32}

\begin{center}
\begin{tikzpicture}[
    every node/.style={font=\ttfamily\small},
    cell/.style={draw,minimum height=0.7cm,align=center},
]
\node[cell,minimum width=4.0cm,fill=mem1!20] (ppn) {ppn};
\node[cell,minimum width=2.0cm,fill=mem2!20,right=0pt of ppn] (flags) {DAGUXWRV};
\node[font=\scriptsize] at ($(ppn.north west)+(0.2,0.2)$) {31};
\node[font=\scriptsize] at ($(ppn.north east)+(-0.3,0.2)$) {20 19};
\node[font=\scriptsize] at ($(flags.north west)+(0.3,0.2)$) {10 9};
\node[font=\scriptsize] at ($(flags.north east)+(-0.3,0.2)$) {flags 0};
\end{tikzpicture}
\end{center}

Los \emph{flags} en una \texttt{pte} son los siguientes bits:

\begin{itemize}[leftmargin=*]
    \item \textbf{V} (\emph{Valid}): indica si la entrada es válida.
    \item \textbf{D} (\emph{Dirty}): encendido por el hardware cuando la
        página correspondiente fue \emph{modificada} (escrita) por una
        instrucción.
    \item \textbf{A} (\emph{Accessed}): encendido por el hardware cuando
        la página fue accedida para lectura o escritura.
    \item \textbf{U} (\emph{User mode}): permiso de acceso en modo
        usuario.
    \item \textbf{R, W, X}: permisos de \emph{Read}, \emph{Write} y
        \emph{eXecute} sobre la página.
    \item \textbf{G} (\emph{Global mapping}): la entrada existe en
        todas las tablas (mapping global).
\end{itemize}

En el proceso de traducción de una dirección lógica a una física, la
CPU verifica los flags correspondientes. Si hay una violación (por
ejemplo, $V=0$, o el acceso no respeta los permisos $R/W/X/U$), se
genera un \emph{page fault}. El kernel atrapa estas excepciones y actúa
en consecuencia.

\paragraph{Context switch.}
El kernel debe cambiar el registro \texttt{satp} en cada context switch
para que apunte a la tabla de páginas de la nueva tarea seleccionada.
Los bits \emph{D} y \emph{A} son la base de muchas políticas de
\emph{memoria virtual} que se ven en el capítulo siguiente.

% =====================================================================
\section{Lookaside Translation Buffers (TLB)}
% =====================================================================

Un sistema con paginado traduce \emph{constantemente} direcciones
lógicas a físicas usando las tablas de páginas. Estas tablas están
almacenadas en RAM y son configuradas por el SO.

\begin{importante}
Sin ninguna optimización, en un esquema multinivel \emph{cada acceso a
memoria del programa requiere varios accesos extra a RAM}: uno por cada
nivel del árbol de tablas (por ejemplo, \emph{pte} raíz, \emph{pte} de
la hoja) y finalmente el acceso al dato. En RV-32 esto multiplicaría
por 3 cada referencia.
\end{importante}

\subsection{La caché de traducciones}

Para evitar esta sobrecarga las CPUs modernas usan una memoria caché
especial dedicada a almacenar las \emph{ptes} más frecuentemente
referenciadas. Esta caché se conoce como
\textbf{Translation Lookaside Buffer} (\textbf{TLB}). La TLB reside
físicamente entre la CPU y la caché de nivel 1 (L1).

\begin{definicion}
La TLB es una memoria \textbf{direccionable por contenido}: se busca
por dirección lógica y, si está presente (\emph{hit}), retorna la
dirección física correspondiente; en caso contrario ocurre un
\emph{miss} y la CPU debe consultar las tablas de páginas en RAM.
\end{definicion}

\subsection{TLB con direct mapping}

Generalmente, una TLB implementa una \textbf{función hash} donde la
clave de búsqueda son una parte de los bits \emph{menos significativos}
del \emph{virtual page number} (vpn), asociados a un \textbf{tag} que
es el conjunto de bits \emph{más significativos} del vpn. El par
(\texttt{tag}, \texttt{index}) identifica unívocamente a un vpn.

\begin{center}
\begin{tikzpicture}[
    every node/.style={font=\ttfamily\small},
    cell/.style={draw,minimum height=0.7cm,align=center},
]
% Virtual address
\node[font=\bfseries\small] at (2.5,2.3) {virtual address (va)};
\node[cell,minimum width=1.6cm,fill=mem1!20] (tag)  at (0.8,1.6) {tag};
\node[cell,minimum width=1.6cm,fill=mem2!20] (idx)  at (2.4,1.6) {index};
\node[cell,minimum width=1.6cm,fill=mem3!20] (off)  at (4.4,1.6) {offset};

\draw[decorate,decoration={brace,amplitude=4pt}]
    ($(tag.south west)+(0,-0.1)$) -- ($(idx.south east)+(0,-0.1)$)
    node[midway,below=4pt,font=\scriptsize\itshape] {virtual page number};

% TLB
\node[font=\bfseries\small] at (8.5,1.6) {TLB entries};
\node[cell,minimum width=4.5cm] (e0) at (8.5,1.0) {\ldots};
\node[cell,minimum width=1.2cm,fill=mem1!20] (et) at (7.0,0.3)  {tag};
\node[cell,minimum width=1.6cm,fill=mem4!20] (ep) at (8.5,0.3)  {ppn};
\node[cell,minimum width=1.2cm] (ef) at (9.9,0.3)  {flags};
\node[cell,minimum width=4.5cm] (e1) at (8.5,-0.4) {\ldots};
\end{tikzpicture}
\end{center}

Cada entrada de la TLB contiene el \texttt{tag}, el \texttt{ppn}
(\emph{physical page number}) y los \texttt{flags}, entre los cuales
está el bit \textbf{valid} (\textbf{V}) que indica si la entrada es
válida.

\paragraph{Condición de \emph{hit}.}
La caché así descripta es \textbf{1-way associative cache} (direct
mapping): hay una sola \texttt{ppn} posible para cada \texttt{index}.
Sea \texttt{va.tag} y \texttt{va.index} las partes de la dirección
virtual recibida; ocurre un \emph{hit} cuando:
\[
\textit{TLB}[\textit{index}].\textit{tag} == \textit{va}.\textit{tag}
\quad\wedge\quad
\textit{TLB}[\textit{index}].\textit{flags}\;\&\;\textit{BIT\_V} \neq 0,
\]
y \texttt{TLB[index].ppn} contiene la \texttt{physical page number}
buscada.

\paragraph{Manejo de un \emph{miss}.}
Si ocurre un \emph{miss}, la CPU accede a la tabla de páginas en RAM
(o a una caché de nivel inferior) y obtiene la \texttt{ppn}
correspondiente. Luego se setea \texttt{TLB[index]} apropiadamente con
los datos de la nueva traducción para acelerar accesos futuros.

\paragraph{Colisiones.}
Inevitablemente, varias direcciones virtuales distintas comparten el
mismo \texttt{index} pero con \texttt{tag} diferente. Eso produce una
\textbf{colisión}: la entrada anterior es reemplazada con el nuevo
trío \texttt{<tag, ppn, flags>}.

\subsection{Otras configuraciones}

Existen organizaciones más sofisticadas que permiten almacenar
\emph{más de una entrada} (un \emph{set}) en el mismo \texttt{index};
son las \emph{set-associative caches} (o \emph{associative caches}). La
configuración \emph{1-way} descripta arriba es el caso más simple. A
mayor asociatividad menor tasa de colisiones, pero mayor complejidad y
costo del hardware de búsqueda.

Algunas arquitecturas, como \textbf{MIPS}, permiten que sea el
\emph{kernel} el que decida (por software) qué hacer ante un
\emph{miss}: el hardware genera una excepción (trap) y el kernel
implementa la política de reemplazo y la carga de la nueva entrada.

% =====================================================================
\section{Cierre y conexión con memoria virtual}
% =====================================================================

A lo largo del capítulo se construye, paso a paso, la maquinaria que
hace posible la gestión moderna de memoria:

\begin{itemize}[leftmargin=*]
    \item La \textbf{jerarquía de memorias} (registros, cachés, RAM,
        disco) explica por qué la velocidad de acceso a la memoria es
        un factor crítico y motiva la existencia de toda la maquinaria
        de caching y traducción.
    \item Las estrategias clásicas de \textbf{asignación} (tamaño
        variable con \emph{first/best/worst fit}, tamaño fijo y
        \emph{buddy system}) muestran el compromiso permanente entre
        \emph{fragmentación interna} y \emph{externa}.
    \item Los mecanismos de \textbf{protección} (registros base/límite,
        segmentación) garantizan que cada proceso quede
        \emph{confinado} a su espacio de memoria y son el origen del
        clásico \emph{segmentation fault}.
    \item El \textbf{código independiente de la posición} (PIC) hace
        que un mismo binario pueda ejecutarse en cualquier dirección
        base, lo que habilita bibliotecas compartidas y técnicas
        modernas como ASLR.
    \item El \textbf{paginado} unifica todo lo anterior: divide la
        memoria física en frames del mismo tamaño y le da a cada
        proceso (y al kernel) un espacio de direcciones lógicas
        independiente. Las tablas de páginas multinivel evitan el
        desperdicio de las tablas planas en espacios de direcciones
        grandes.
    \item Las \textbf{TLBs} aceleran la traducción dirección
        lógica $\to$ dirección física, evitando que cada acceso a
        memoria pague el costo de recorrer varios niveles de tablas en
        RAM.
\end{itemize}

\begin{idea}
Esta maquinaria es exactamente la base sobre la que se construyen las
políticas de \textbf{memoria virtual} (swapping, demand paging,
copy-on-write, mapeo de archivos, \ldots) que se estudian en el
capítulo siguiente. Los bits \texttt{V}, \texttt{D} y \texttt{A} de la
\emph{pte}, y la posibilidad de generar \emph{page faults}, son las
herramientas que el SO necesita para implementar esas políticas.
\end{idea}

\end{document}
